%قبل از شروع مسئله ترجیحاً چیزی ننویسید. اما اگر خواستید بنویسید، در قالب لیستی (فقرات و شمارش) بنویسید.
%برای اضافه کردن نمره، بعدی از شروع‌مسئله، [x] را اضافه کنید و به جای x  نمره مورد نظر را قرار دهید. برای زیرمسئله‌ها هم به همین صورت است.
%\شروع‌مسئله[10]
\شروع‌مسئله [20]
% صورت اصلی مسئله را در این‌جا بنویسید.
فرض کنید $H$ یک خانواده از توابع درهم‌سازی باشد که هر تابع درهم‌سازی $h \in H$ مجموعۀ $U$ از کلیدها را به $\{0, 1, ..., m-1\}$ می‌برد. می‌گوییم خانوادۀ $H$ از توابع درهم‌سازی \متن‌لاتین{k-universal} است اگر برای هر $k$تایی متمایز $(x_1, x_2, ..., x_k) \in U^k$ و هر $k$ تایی $(y_1, y_2, ..., y_k) \in [m]^k$ و هر $h \in H$ که به صورت تصادفی انتخاب شده باشد داشته باشیم:
$$Pr(h(x_1) = y_1 \wedge h(x_2) = y_2 \wedge ...\wedge h(x_k) = y_k) = m^{-k}.$$
به عبارت دیگر، $k$ تایی $(h(x_1), h(x_2), ..., h(x_n))$ با احتمالی برابر می‌تواند هر کدام از $k$ تایی‌های مرتب از مجموعه‌ی $[m]^k$ باشد.

\شروع{زیرمسائل}

\زیرمسئله [4] ثابت کنید اگر $H$ یک خانوادۀ \متن‌لاتین{2-universal} باشد، \متن‌لاتین{universal} نیز هست.

\زیرمسئله [6] فرض کنید $U$ مجموعۀ تمام $n$تایی‌ها از $Z_p = \{0, 1, ..., p-1\}$ باشد که $p$ یک عدد اول است. کلید $x = (x_1, x_2, ..., x_n) \in U$ را در نظر بگیرید. برای هر $n$ تایی $a = (a_1, a_2, ..., a_n) \in U$ تابع درهم‌سازی $h_a$ را به این صورت تعریف می‌کنیم
$$h_a(x) = (\sum_{j=1}^{n} a_jx_j)\: mod\: p$$
فرض کنید $H$ خانوادۀ تمام $h_a$ ها باشد. ثابت کنید $H$ یک خانوادۀ \متن‌لاتین{universal} است ولی \متن‌لاتین{2-universal} نیست.

\زیرمسئله [6] حال خانوادۀ $H$ در قسمت قبل را کمی تغییر می‌دهیم. برای هر $a \in U$ و هر $b \in Z_p$ تابع $h^\prime_{ab}$ را به این صورت تعریف می‌کنیم
$$h^\prime_{ab}(x) = (\sum_{j = 1}^{n} a_jx_j + b)\: mod\: p$$
فرض کنید $H^\prime$ خانواده‌ی تمام $h^\prime_{ab}$ ها باشد. ثابت کنید $H^\prime$ یک خانوادۀ \متن‌لاتین{2-universal} است.

\زیرمسئله [4] فرض کنید علی و رضا بین خودشان یک تابع درهم‌سازی $h$ را از خانوادۀ \متن‌لاتین{2-universal} که
	$H$ 
	نام دارد
	انتخاب کنند و آن را با هیچ کس در میان نگذارند. هر تابع درهم‌سازی $h \in H$ مجموعۀ $U$ از کلیدها را به $Z_p$ می‌برد که $p$ یک عدد اول است. علی از طریق اینترنت برای رضا پیام $m \in U$ را می‌فرستد. او برای اطمینان از درست رسیدن پیامش $t = h(m)$ را نیز همراه پیامش می‌فرستد. رضا بعد از دریافت پیام چک می‌کند که آیا رابطه‌ی $t = h(m)$ برقرار است یا نه. فرض کنید فردی در مسیر فرستاده شدن پیام مداخله کند تا پیامی اشتباه به دست رضا برسد. او برای این کار $(m, t)$ را با $(m^\prime, t^\prime)$ عوض می‌کند تا رضا را فریب دهد. نشان دهید حتی در صورتی که آن فرد خانواده‌ی $H$ را بداند، احتمال این که بتواند رضا را فریب دهد طوری که رضا $(m^\prime, t^\prime)$ را قبول کند حداکثر $\frac{1}{p}$ است.




\پایان{زیرمسائل}
%اگر مسئله چند قسمت داشت به این صورت عمل کنید:
%\شروع{زیرمسائل}
%\زیرمسئله[0] % مسئلۀ ۱آ
%\زیرمسئله % مسئلۀ ۱ب
%\زیرمسئله % مسئلۀ ۱پ
%\زیرمسئله % مسئلۀ ۱ت
%\پایان{زیرمسائل}

\solution{
	
آ. اگر $H$ یک خانوادۀ \متن‌لاتین{2-universal} باشد آنگاه برای هر دو کلید $x \neq y$ زوج $(h(x), h(y))$ با احتمال یکسانی می‌تواند هر کدام از اعضای مجموعۀ $\{0, 1, ..., m-1\} \times \{0, 1, ..., m-1\}$ باشد. در $m$ حالت از $m^2$ حالتی که این زوج می‌تواند داشته باشد برخورد داریم. پس اگر تمام $h$ها را در نظر بگیریم در $\frac{1}{m} |H|$ از آن‌ها برخورد رخ می‌دهد پس $H$ یک خانواده‌ی \متن‌لاتین{universal} است.

ب. دو کلید متفاوت $x$ و $y$ را در نظر بگیرید. بدون کاسته شدن از کلیت مسأله فرض کنید این دو کلید در جایگاه اول متفاوت باشند. برخورد در صورتی رخ می‌دهد که داشته باشیم:
$$a_1(x_1 - y_1) + a_2(x_2 - y_2) + ... + a_{n}(x_{n} - y_{n}) \equiv 0 \ \ (mod \ \ p)$$
$$\Leftrightarrow a_1(x_1 - y_1) \equiv -(a_2(x_2 - y_2) + ... + a_{n}(x_{n} - y_{n}) ) \ \ (mod \ \ p)$$
حال $a_2$ تا $a_n$ را ثابت نگاه می‌داریم و $a_1$ را تغییر می‌دهیم. دقت کنید سمت راست عددی ثابت است. همچنین می‌دانیم که باقیماندۀ $x_1 - y_1$ به  $p$ برابر ۰ نیست. پس از $p$ حالتی که می‌توانیم $a_1$ را مقدارگذاری کنیم دقیقاً یکی وجود دارد که رابطۀ بالا برقرار شود. پس حداکثر برای $\frac{1}{p}$ از اعضای $H$ داریم $h_a(x) = h_a(y)$. پس این خانواده \متن‌لاتین{universal} است.
حال کلید $x = (0, 0, ..., 0)$ و یک $y$ ثابت را در نظر بگیرید. برای هر عضو از $H$ داریم $(h_a(x), h_a(y)) = (0, h_a(y))$. پس $(h_a(x), h_a(y)$ نمی‌تواند با احتمال یکسانی برابر هر کدام از اعضای $\{0, 1, ..., p-1\} \times \{0, 1, ..., p-1\}$ باشد. پس این خانواده \متن‌لاتین{2-universal} نیست.

پ. دو کلید متفاوت $x$ و $y$ را در نظر بگیرید و مانند قبل فرض کنید در جایگاه اول متفاوت هستند. $a_2$ تا $a_n$ را ثابت می‌گیریم. ثابت می‌کنیم برای هر $r$ و $t$ از $\{0, 1, ..., p-1\}$ دقیقا یک زوج $(a_1, b)$ وجود دارد به طوری که $(h(x), h(y)) = (r, t)$. باید داشته باشیم:
$$r \equiv a_1x_1 + b + X \ \ (mod \ \ p)$$
$$t \equiv a_1y_1 + b + Y \ \ (mod \ \ p)$$
$$\Leftrightarrow$$
$$a_1 \equiv (r - t - X + Y)(x_1 - y_1)^{-1} \ \ (mod\ \ p)$$
$$b \equiv r - a_1x_1 - X \ \ (mod \ \ p)$$
طبق فرض می‌دانیم که $(x_1 - y_1)^{-1}$ به صورت یکتا وجود دارد. پس حکمی که می‌خواستیم ثابت کنیم ثابت شد. پس برای هر انتخاب $a_2$ تا $a_n$ هر زوج $(r, t)$ دقیقا یک بار تولید می‌شود. پس هر زوج $(r, t)$ دقیقا $p^{n-1}$ بار تولید می‌شود. پس این خانواده \متن‌لاتین{2-universal} است.

ت. برای کلیدهای $m$ و $m^{\prime}$ اگر $h$ تصادفی انتخاب شود، زوج $(h(m), h(m^{\prime}))$ با احتمال یکسانی می‌تواند هر کدام از اعضای $\{0, 1, ..., p-1\} \times \{0, 1, ..., p-1\}$ باشد. پس تمام زوج‌های $(t, h(m^{\prime}))$ با احتمال یکسان $\frac{1}{p}$ می‌توانند رخ دهند. پس فردی  که می‌خواهد پیام را عوض کند با دانستن $H$ و با دیدن $(m, t)$ هیچ چیزی درباره‌ی $h(m^\prime)$ نمی‌فهمد. او به دنبال زوج $(t, h(m^{\prime}))$ ای است که احتمال رخ دادن آن بیشینه باشد. از آن جایی که تمام این زوج‌ها با احتمال یکسان $\frac{1}{p}$ رخ می‌دهند، پس او باید یکی از آن‌ها را تصادفی انتخاب کند. در نتیجه احتمال اینکه رضا پیام را قبول کند برابر $\frac{1}{p}$ خواهد بود.

}




%%%%%%%%%%%%%%%%%%%%%%%%%%%%%%%%%%%%%%%%%%%%%%%%%%%%%
% برخی دستورات مفید برای تمرین‌ها. %
%%%%%%%%%%%%%%%%%%%%%%%%%%%%%%%%%%%%%%%%%%%%%%%%%%%%%

% فرمول‌های نمونه:
%$f(x) = \Theta(x)$
%$T(x, y) \leq \log(x) + 2^y + \binom{2n}{n}$
% {\tt code\_function}


% لیست موارد بدون شماره‌گذاری:
%\شروع{فقرات}
%    \فقره مورد اول 
%        \شروع{فقرات}
%            \فقره مورد اول داخلی 
%                \شروع{فقرات}
%                    \فقره اولین مورد خیلی داخلی 
%                    \فقره دومین مورد خیلی داخلی 
%                \پایان{فقرات}
%            \فقره مورد دوم داخلی 
%        \پایان{فقرات}
%    \فقره مورد دوم 
%\پایان{فقرات}

% لیست موارد شماره‌گذاری شده:
%\شروع{شمارش}
%    \فقره یک مورد 
%    \فقره یک مورد دیگر 
%    \فقره باز هم یک مورد دیگر
%\پایان{شمارش}

% برای نوشتن یک فرمول چند خطی که یک شماره بخوره. برای حالت‌های دیگر به سایت زیر مراجعه کنید:
% https://www.overleaf.com/learn/latex/Aligning_equations_with_amsmath
%\begin{align} 
%    \begin{split}
%        (x+y)^3 &= (x+y)^2(x+y) \\
%                &= (x^2+2xy+y^2)(x+y) \\
%                &= (x^3+2x^2y+xy^2) + (x^2y+2xy^2+y^3) \\
%                &= x^3+3x^2y+3xy^2+y^3
%    \end{split}                                 
%\end{align}

% نوشتن ماتریس و آرایه دوبعدی
%\begin{align}
%    A = 
%    \begin{bmatrix}
%        A_{11} & A_{21} \\
%        A_{21} & A_{22}
%    \end{bmatrix}
%\end{align}

% اضافه کردن تصویر با شرح. برای اضافه کردن تصویر بدون شرح، فقط دستور درج‌تصویر را استفاده کنید.
%\شروع{شکل}[h!]
%\تنظیم‌ازوسط
%  \درج‌تصویر[پهنا=0.5\پهنای‌متن]{img.jpg}
%\شرح{شرح تصویر}
%\برچسب{تصویر:۱}
%\پایان{شکل}
%\پایان{زیرمسائل}

%کد زیر را ببینید.
%%\lstinputlisting{source_filename.py}
%\begin{latin}
%\begin{lstlisting}
%# Import the modules
%import sys
%import random
%
%ans = True
%
%while ans:
%question = raw_input("Ask the magic 8 ball a question: (press enter to quit) ")
%
%answers = random.randint(1,8)
%
%if question == "":
%sys.exit()
%
%elif answers == 1:
%print "It is certain"
%
%elif answers == 2:
%print "Outlook good"
%
%elif answers == 3:
%print "You may rely on it"
%
%elif answers == 4:
%print "Ask again later"
%
%elif answers == 5:
%print "Concentrate and ask again"
%
%elif answers == 6:
%print "Reply hazy, try again"
%
%elif answers == 7:
%print "My reply is no"
%
%elif answers == 8:
%print "My sources say no"
%\end{lstlisting}
%\end{latin}
